%% The following is a directive for TeXShop to indicate the main file
%%!TEX root = diss.tex
%==============================================================================
\chapter{Further research directions}
%==============================================================================

In this thesis we explored the usefulness of using operators to simplify the derivation of functional equations pertaining to random walk models on $\ze^d$ and to simplify the first steps in solutions to the problems. The sleek approach of operators makes us wonder if we can extend the work further into territories like Sturm-Liouville theory where we may investigate the spectra of operators and find alternative solutions to the functional equations via eigenfunctions and eigenvalues of these operators. For this to work we need a suitable inner product for which we may investigate if the operators are self-adjoint and so on. Unfortunately, we haven't found such suitable inner products apart from inner products on symmetric functions which may be promising. In this chapter we'll only provide the start of such analysis which may be undertaken as future research.

%=======================================
\section{Operators}
%=======================================

The operators we've seen in this thesis are almost invariably marking operators (like $L$ and $E$), marking operators in combination of multiplication by another function (like $\Gamma S$) or a combination of operators (like $L^{\um 1} - tES$). In this order we'll investigate the operators' properties, eigenvalues, eigenvectors, kernels and invertibility.


%=======================================
\subsection{Marking operators}
%=======================================

The building block of the operators are coefficient extractions $[x_I^J]$ and in our case we see term extractions $\ext{x_I^J}$. We'll use the notation
\begin{definition} For sets $I$ and $J$ of equal size we define the term extraction operator as
\begin{align}
	[[x_I^J]] \equiv \ext{x_I^J}\,.
\end{align}
\end{definition}
The main question here is how to compose these term extractions. Suppose that we want to perform the extractions with monomials $x_I^P$ and $x_J^Q$. Then the composition will contain three extractions - one for the $x$ powers appearing in $x_I^P$ but not in $x_J^Q$, one for the reversed situation and one for the common terms between the two monomials. However the resulting composition will be the zero operator if the power sequences for the common variables don't match.
\begin{example} We compute
\begin{align}
	[[x_{1,2}^{3,4}]][[x_{2,3}^{4,1}]] = [[x_{1,2,3}^{3,4,1}]]
\end{align}
and
\begin{align}
	[[x_{1,2}^{3,2}]][[x_{2,3}^{4,1}]] = 0\,.
\end{align}
\end{example}
The next highly-technical lemma makes the composition precise.
\begin{lemma} Let $I,\, J,\, P,\, Q$ be sets and $K = I\cap J$, $I' = I\backslash K$, $J' = J \backslash K$. Then let $P' = P|_{I'}$, $Q' = Q|_{J'}$ and $P'' = P\backslash P'$, $Q'' = Q\backslash Q'$. Then the composition of two term extractions is given by
\begin{align}
	[[x_I^P]][[x_J^Q]] = \delta(P'',Q'')[[x_{I'}^{P'}x_{J'}^{Q'}x_K^{P''}]]
\end{align}
That is, the resulting operator is zero if the power sequences $P$ and $Q$ don't match over the variables common to the two extractions. Furthermore the two term extraction operators commute.
\end{lemma}

Combinatorially speaking, it should be clear that if a marking operator marks several different monomials we may express this operator as a product of marking operators each marking a single monomial. We can make this notion precise with the algebraic form of the marking operators.
\begin{definition} The operator $\Gamma$ that marks a monomial $x_I^J$ with a label $a$ is given by
\begin{align}
	\Gamma = 1  +(a-1)\ext{x_I^J} = 1 + b[[x_I^J]]\,.
\end{align}
A finite marking operator is one marking finitely many monomials. An infinite marking operator is one that is not finite.
\end{definition}
\begin{lemma} Let $\alpha$ and $\beta$ be the labels for two different monomials $x_I^P$ and $x_J^Q$ that are marked by $\Gamma$. Then
\begin{align}
	\Gamma = 1 + (\alpha-1)[[x_I^P]]+(\beta-1)[[x_J^Q]] = (1+ (\alpha-1)[[x_I^P]])(1+ (\beta-1)[[x_J^Q]])\,.
\end{align}
\end{lemma}
\begin{corollary} Marking operators commute.
\end{corollary}
\begin{proof} This can be seen from their combinatorial nature. Since multiplication is commutative and the marking is done via multiplication, the order of applying the marking operators doesn't matter. Algebraically this is a consequence of the commutativity the term extraction operators.
\end{proof}

Combinatorially speaking, we can find some of the eigenvalues and eigenfunctions easily by looking at single monomials. Every non-marked monomial is an eigenfunction with eigenvalue 1. Every marked monomial is an eigenfunction with eigenvalue equal to the label. If we seek eigenfunctions among polynomials of the $x$ variables by linear independence of the indeterminates we'll have that the eigenpairs listed so far are all we have. However this doesn't say anything about eigenvalues for eigenfunctions that are power series or Laurent series in the variables $x$. Although the analysis is straight-forward for operators marking single monomials.
\begin{prop} Let $\Gamma$ be an operator marking a single monomial $x_I^J$ with label $a$. Let $\lambda$ and $u_\lambda$ be an eigenvalue and $u$ be an eigenfunction corresponding to that eigenvalue. Then $\lambda = a$ with $u_\lambda = x_I^J$ or $\lambda = 1$ with $u_\lambda$ a function not containing $x_I^J$ (in the formal power series sense). In particular every monomial other than $x_I^J$ is an eigenfunction with eigenvalue 1.
\end{prop}
\begin{proof}
Let $u_\lambda = u' + u''x_I^J$ where $u \neq 0$ and $u''$ doesn't share variables with $x_I^J$. Then
\begin{align}
	\Gamma(u) = u' + au''x_I^J  = \lambda (u' + u''x_I^J)\,,
\end{align}
which we may write as
\begin{align}
	(1-\lambda) u' + (a-\lambda)u''x_I^J  = 0\,.
\end{align}
We have two cases.
\begin{itemize}
\item If $u' = 0$ then $u'' \neq 0$ and $\lambda = a$. The eigenfunction is $x_I^J$.
\item If $u' \neq 0$ and there is no term in $u''x_I^J$ that can compesate for this. Therefore $\lambda = 1$ and $u'' = 0$.
\end{itemize}
\end{proof}

\begin{corollary} The kernel of a finite marking operator $\Gamma$ is generated by the monomials marked with label 0.
\end{corollary}

The proofs can be extended for finite marking operators. It's not entirely clear how we may prove a similar result for operators marking infinitely many monomials.

%=======================================
\subsection{Marking operators with multiplication}
%=======================================
 We'll call an operator inhomogeneous if it consists of a homogeneous operator $\Gamma$ and a laurent polynomial (possibly infinite) $S$ and and acts in one of the following ways
 \begin{align}
	\Gamma S(u) = \Gamma(Su)\quad\text{or}\quad S\Gamma u = S\cdot \Gamma u\,.
 \end{align}
This kind of operator is important because it can be found in the functional equations for random walk models. We seek to find the eigenpairs for  inhomogeneous operators.  We first show that the two definitions are equivalent when it comes to finding eigenpairs.
\begin{prop} The two operators have the same eigenvalues with eigenfunctions differing by a multiple of $S$. That is $(\lambda,u)$ is an eigenpair for $\Gamma S$ if and only of $(\lambda, Su)$ is an eigenpair for $S \Gamma$. In particular, if $\lambda = 0$ we have
\begin{align}
	\ker(\Gamma S) = \ol{S} \ker(\Gamma)\,.
\end{align}
\end{prop}
\begin{proof} Note that
\begin{align}
	\Gamma Su = \lambda u \quad \Longleftrightarrow \quad S\Gamma (Su) = \lambda (Su)\,.
\end{align}
Therefore the operators have the same eigenvalues and they have the eigenfunctions $u$ and $Su$ respectively. Now if we pick $\lambda = 0$ to find the kernels of the operators we find that $u \in \ker(\Gamma S)$ if and only if $Su \in \ker(S\Gamma) = \ker(\Gamma)$ which completes the proof.
\end{proof}
It is a rather difficult task to find the eigenpairs for such operators as the next example shows
\begin{example}
Let $\Gamma$ and an eigenfunction $u$ of $\Gamma S$ where 
\begin{align}
	S = x\com u = \sum_k u_k x^k\,.
\end{align}
Computing $(\Gamma S-\lambda I)u$ leads to the system of equations
\begin{align}
	\lambda u_n &= a u_{n-1}\,,\\
	\lambda u_k &= u_{k-1}, \quad \text{if } k \neq n\,.
\end{align}
We have a doubly-infinite recursion so we can pick an arbitrary initial condition. If $u_n = 0$ the rest of the terms will be zero as well so we must have $u_n \neq 0$. Set $u_n = 1$. Then via induction we can find
\begin{align}
	u_{n+k} &= 1/\lambda^{k}, \quad \text{for $k \ge 0$}\,,\\
	u_{n-k} &= \lambda^{k}/a^k, \quad \text{for $k \ge 1$}\,,
\end{align}
and this shows that we have an eigenfunction for each $\lambda \neq 0$. Computing $u(x)$ formally leads to
\begin{align}
	u(x) &=-\frac{\lambda b x^{n+1}}{(\lambda-x)(\lambda-ax)}\,.
\end{align}
\end{example}
The situation becomes more complicated for more complicated functions $S$.
%=======================================
\subsection{Random Walks Operator}
%=======================================
In this section we explore the operator $\mc{L} = 1 - t\Gamma S$ for a step Laurent polynomial $S$ and some random walk operator $\Gamma$. Let $(u, \lambda)$ be an eigenpair for $\mc{L}$. Then we calculate
\begin{align}\begin{split}
	(\mc{L}-\lambda I) u &= u - \lambda u -t\Gamma Su\\
	t \Gamma S u &= (1-\lambda)u = \mu u\,.
\end{split}\end{align}
If $\lambda = 1$ we would have to compute the kernel of $\Gamma S$ which we did in the previous section - the eigenfunctions of $\mc{L}$ for eigenvalue $\lambda = 1$ are the functions $u \in \ol{S}\ker{\Gamma}$. Let $u$ be given by a finite order Laurent series
\begin{align}
	u = u_n t^n + u_{n+1}t^{n+1} + \ldots
\end{align}
where $u_n \neq 0$. Then we can re-write the functional equation as
\begin{align}
	t^{-n}(1-t\Gamma S)u = \lambda t^{-n}u\,.
\end{align}
Letting $t \to 0$ in this equation yields
\begin{align}
	(\lambda - 1)u_n = -\lim_{t\to 0} t^{1-n}\Gamma S u = -\Gamma S \lim_{t\to 0} (t^{1-n} u) = 0\,.
\end{align}
Since $u_n \neq 0$ this implies that $\lambda = 1$ and again we see that $u \in \ol{S}\ker \Gamma$. Since $S$ is a Laurent polynomial, the resulting functions will indeed be finite order Laurent series. This analysis shows that the only eigenfunctions of $\mc{L}$ among finite order Laurent series in $t$ are those that are members of $\ol S \ker \Gamma$ with eigenvalue $\lambda = 1$.
\begin{example} Let $\Gamma = LE$ in the one-dimensional problem with small steps on the nonnegative line with nontrivial operator $E$. Then $\ker \Gamma = \ker E$ consists of functions $\alpha(t) \ol x$ and the eigenfunctions of $\mc{L}$ are $\alpha(t)\ol x \ol S$ for various $\alpha(t)$. In our case we only take finite order Laurent series of $\alpha(t)$. Note that
\begin{align}
	Q = \frac{t(\ol x - \ol \sigma)}{(1-tS)(b-a\ol \sigma t)} = \frac{1}{xS}\frac{\sigma - x}{b\sigma - at}\frac{tS}{1-tS}
\end{align}
where we may write
\begin{align}
	\frac{\sigma - x}{b\sigma - at} = \alpha(t) + \beta(t)x
\end{align}
where $\alpha(t)$ and $\beta(t)$ are algebraic functions of $t$ and $a$ since $\sigma = \sigma(t)$. This at least suggests the possibility of connecting the eigenfunctions of $\mc L$ to the solution $Q$
\end{example}

If we are to have a theory similar to Sturm-Liouville theory we still need to find a suitable inner product on multivariate Laurent series that will allow us to write solutions for $Q$ in terms of those eigenfunctions with eigenvalue 1.


\endinput
%=======================================
%=======================================
%=======================================

%================================
\section{Connection with Fredholm Theory}
%================================
Suppose that we are given the equation in the form
\begin{align}
	KQ +t\Gamma'SQ = q
\end{align}
where
\begin{align}
	\Gamma' = \sum_{I,J}a_{I,J} x_I^J[x_I^J]
\end{align}
where we only have finite-sized $J$'s and $\emptyset \neq I\subseteq [d]$. From complex analysis we know that for a function $f(x)$ given by a power series and for $x$ within the region of convergence
\begin{align}
	[x^n]f(x) = f_n = \frac{1}{2\pi i}\int_\gamma \frac{f(z)}{z^{n+1}}dz
\end{align}
and
\begin{align}
	f(x) = \frac{1}{2\pi i}\int_\gamma \frac{f(z)}{z-x}dz
\end{align}
where in the first case $\gamma$ winds CCW around the origin within the region of convergence and in the second case $\gamma$ winds around $x$ within the region of convergence. The idea is to combine these two expressions for many dimensions (at least formally) to represent the operator $\mc{L}$ in integral form. We have for $|I| = n$
\begin{align}
	[x_I^J]f(x_1,\ldots,x_d)= \frac{1}{(2\pi i)^{n}}\int\cdots\int_\gamma \frac{f}{z_I^{J+1}}dz_I
\end{align}
where $f$ has $x_i = z_i$ for $i \in I$ and the rest are still $x$'s. If we use this representation we'll have different numbers of iterations of the integral for different coefficient. We might as well go all-out and write
\begin{align}
	[x_I^J]f(x_1,\ldots,x_d)= \frac{1}{(2\pi i)^d}\int\cdots\int_\gamma \frac{f}{z_I^{J+1}(z-x)_{I^c}}dz_{[d]}
\end{align}
where now $x_i = z_i$ in $f$ and 
\begin{align}
	(z-x)_P = \prod_{p\in P} (z_p-x_p)
\end{align}
This allows us to write The operator like this
\begin{align}
	(K+t\Gamma'S)Q&= \frac{1}{(2\pi i)^d}\int\cdots\int_\gamma \left(\frac{K}{(z-x)_{[d]}} +t S\sum_{I,J}\frac{a_{I,J}}{z_I^{J+1}(z-x)_{I^c}}\right)Q dz_{[d]}\\
	& = \frac{1}{(2\pi i)^d}\int\cdots\int_\gamma \frac{1}{(z-x)_{[d]}}\left(K +t S\sum_{I,J}\frac{a_{I,J}(z-x)_I}{z_I^{J+1}}\right)Q dz_{[d]}\\
	&= \frac{1}{(2\pi i)^d}\int\cdots\int_\gamma R_0Q dz_{[d]}
\end{align}
Now to turn this into a $d$-dimensional Fredholm problem we'll need to parametrize the integral. One way that might work is to set $z_k = r_k \exp(i\theta_k)$ and $x_k = s_k e^{i\phi_k}$. Since $dz_k = r_k i e^{i\theta_k}$ we have
\begin{align}
(K+t\Gamma'S)Q = \frac{r_{[d]}}{(2\pi)^d}\int_0^{2\pi}\cdots\int_0^{2\pi} e^{i\theta_1}\ldots e^{i\theta_d}R_0Q d\theta_{[d]} = q(\phi_1,\ldots,\phi_d)
\end{align}
so if we write
\begin{align}
R = \exp(i(\theta_1+\ldots + \theta_d))R_0(\phi_1,\ldots,\phi_d,\theta_1,\ldots,\theta_d)
\end{align}
we have the inhomogeneous Fredholm equation of the first kind in $d$ dimensions
\begin{align}
	\int_0^{2\pi}\cdots\int_0^{2\pi} RQ d\theta_1\ldots d\theta_d= g
\end{align}
where $R$ is the kernel and $g = (2\pi)^dq/r_{[d]}$ and $Q$ is in terms of $\theta_k$'s. Perhaps Fredholm theory can be applied in solving the random walk models in many dimensions. Perhaps this will also lead to a good inner product on the space where $Q$ lives.